\subsection{Trajectory advantages, rollout buffers, and policy versions}

The reward in outcome-based language-model RL normally evaluates a complete sampled
response, not an individual token. For one prompt, GRPO samples a group of trajectories,
scores each trajectory, and normalizes its reward relative to the group. The resulting
advantage $A_i$ is therefore one scalar attached to trajectory $i$. It says whether that
response was better or worse than the other responses to the same prompt; it does not say
which token caused the outcome \citep{guo2025deepseekr1}.

To see how this scalar changes the model, first omit clipping and KL control. A simplified
trajectory objective is
\begin{equation}
  J_i(\theta)
  = A_i\frac{1}{T_i}\sum_{t=1}^{T_i}
    \log \pi_\theta(o_{i,t}\mid q,o_{i,<t}),
  \qquad \mathcal{L}_i=-J_i.
  \label{eq:trajectory-advantage-objective}
\end{equation}
The sampled trajectory is replayed through the trainable model with teacher forcing to
recover the log-probability of each sampled action. The trajectory advantage is broadcast
over those actions. If $A_i>0$, minimizing the negative objective increases their
log-probabilities; if $A_i<0$, it decreases them. The batch reduction produces one scalar
loss, and backpropagation adds all token contributions in the shared Transformer
parameters. Thus a successful response reinforces every trainable action in that response,
including actions that may not have caused success. This is outcome-level rather than
process-level credit assignment.

\paragraph{What the rollout buffer records.}
Generation, reward evaluation, and optimization are separate phases. The generated token
IDs, behavior-policy log-probabilities, masks, rewards, and advantages are held in a
short-lived rollout buffer. The notation $\pi_{\mathrm{old}}$ identifies the policy that
generated those stored actions. It is a provenance record for the current buffer, not a
recommendation to generate future trajectories with an inferior model. The current
trainable policy is $\pi_\theta$.

\paragraph{Where the probability ratio comes from.}
The ratio is not a second heuristic signal added to the advantage. It follows from asking
for an expectation under the new policy while holding samples from the old policy. Fix a
prompt-prefix state $s$ and treat the buffered advantage
$A_{\mathrm{old}}(s,a)$ as constant. The new policy's expected advantage is
\begin{align}
  J_s(\theta)
    &= \sum_a \pi_\theta(a\mid s)A_{\mathrm{old}}(s,a) \notag\\
    &= \sum_a \pi_{\mathrm{old}}(a\mid s)
       \frac{\pi_\theta(a\mid s)}{\pi_{\mathrm{old}}(a\mid s)}
       A_{\mathrm{old}}(s,a) \notag\\
    &= \mathbb{E}_{a\sim\pi_{\mathrm{old}}}
       \left[\rho_\theta(s,a)A_{\mathrm{old}}(s,a)\right].
  \label{eq:importance-weighted-expected-advantage}
\end{align}
This change of sampling distribution is valid where the old policy gives the action
nonzero probability. The advantage supplies desirability and scale. The ratio
$\rho_\theta=\pi_\theta/\pi_{\mathrm{old}}$ supplies no new preference; it reweights each
old-policy sample according to how much probability mass the new policy assigns it.

For example, suppose the old policy chooses a good action with probability $0.25$ and a
bad action with probability $0.75$, with advantages $+3$ and $-1$. Its expected advantage
is zero. If the new probabilities are $0.5$ and $0.5$, its expected advantage is one.
Reweighting old-policy samples gives the same value:
$0.25(0.5/0.25)(3)+0.75(0.5/0.75)(-1)=1$. Omitting the ratios would still compute the old
expectation, zero, rather than the new one.

The advantage by itself cannot update the model because it is detached:
$\nabla_\theta A_{\mathrm{old}}=0$. Parameter dependence enters through
$\log\pi_\theta$ in Equation~\ref{eq:trajectory-advantage-objective}, or through
$\rho_\theta$ in Equation~\ref{eq:importance-weighted-expected-advantage}. At
$\pi_\theta=\pi_{\mathrm{old}}$, these forms have the same first-order policy gradient
because $\nabla_\theta\rho=\rho\nabla_\theta\log\pi_\theta$. The canonical PPO/GRPO
surrogate is therefore $\rho A$; it does not require multiplying $\rho A$ by a separate
cross-entropy term.

For a complete response, exact importance sampling would use the trajectory ratio
\begin{equation}
  \frac{\pi_\theta(o\mid q)}{\pi_{\mathrm{old}}(o\mid q)}
  = \prod_{t=1}^{T}
    \frac{\pi_\theta(o_t\mid q,o_{<t})}
         {\pi_{\mathrm{old}}(o_t\mid q,o_{<t})}.
  \label{eq:complete-trajectory-importance-ratio}
\end{equation}
This product becomes high-variance for long responses. PPO-style training instead keeps
the sampled prefix distribution fixed and constructs a local per-token surrogate with
$\rho_{i,t}A_i$. GRPO broadcasts its outcome advantage across the trajectory. This is a
first-order improvement objective around $\pi_{\mathrm{old}}$, not exact reweighting of a
distant new-policy trajectory distribution; clipping helps enforce that local use
\citep{schulman2017ppo}.

\paragraph{The same derivation under asynchronous policy lag.}
Asynchronous RL changes the source of the distribution mismatch, not the importance-
weighting mathematics. Suppose a rollout worker downloads policy version 10 and starts a
long tool-using trajectory. While it generates, the learner advances to version 13 using
other workers' data. When the trajectory arrives, its behavior policy is
$\pi_{\mathrm{old}}=\pi_{10}$ and its exact sampled tokens and behavior log-probabilities
must remain attached to it. The learner evaluates those actions under
$\pi_\theta=\pi_{13}$ and forms
\begin{equation}
  \rho_{i,t}
  = \frac{\pi_{13}(o_{i,t}\mid q,o_{i,<t})}
         {\pi_{10}(o_{i,t}\mid q,o_{i,<t})}.
  \label{eq:asynchronous-policy-lag-ratio}
\end{equation}
This is Equation~\ref{eq:importance-weighted-expected-advantage} with an older worker as
the sampling distribution. In synchronous reuse, the policies match during rollout and
diverge when learner steps reuse the buffer. In asynchronous RL, they may already differ
when the trajectory reaches the learner because generation and queueing took time.

Importance weighting does not fully rescue arbitrarily stale data. Exact correction of the
prefix distribution would require the high-variance trajectory product in
Equation~\ref{eq:complete-trajectory-importance-ratio}; the per-token surrogate only makes
a local correction on prefixes sampled by the old worker. Large policy lag produces
extreme ratios, extensive clipping, and little reliable improvement signal. Consequently,
asynchronous systems must combine ratios with weight refreshes, queue controls, policy-lag
limits, and filtering or rejection of stale trajectories. Re-tokenizing a response or
recomputing its denominator with the learner would erase the behavior-policy provenance
and invalidate the ratio.

At the start of a synchronous rollout iteration, the policies are equal. The system sets
$\pi_{\mathrm{old}}\leftarrow\pi_\theta$ and samples the buffer. If it evaluates one
full-batch loss, takes exactly one optimizer step, discards the buffer, and immediately
resamples, then the importance ratio
\begin{equation}
  \rho_{i,t}(\theta)
  = \frac{\pi_\theta(o_{i,t}\mid q,o_{i,<t})}
         {\pi_{\mathrm{old}}(o_{i,t}\mid q,o_{i,<t})}
  \label{eq:rollout-buffer-importance-ratio}
\end{equation}
equals one when the loss is evaluated. Clipping is inactive in this schedule. The gradient
is nevertheless nonzero because the old probability is detached whereas the current
probability is differentiable: at $\rho=1$,
$\nabla_\theta(\rho A_i)=A_i\nabla_\theta\log\pi_\theta$.

\paragraph{Why reuse a buffer at all?}
Reuse is an optional compute tradeoff, not a requirement of the advantage calculation.
Autoregressive LLM rollout generation requires a sequential forward pass for every token
and may also wait for tools, environments, or verifiers. Once a system has paid that cost,
it may perform several minibatch steps or optimization epochs on the same buffer. Several
small steps can move the policy farther than one conservative step without attempting one
large update from a noisy Monte Carlo batch. Even a single pass over a large buffer may use
many optimizer steps: after the first minibatch, $\pi_\theta$ has changed, while every
remaining minibatch still contains actions from $\pi_{\mathrm{old}}$.

This reuse creates the mismatch that the PPO-style ratio and clipping control
\citep{schulman2017ppo}. It does not create new information. Excessive reuse can overfit the
sampled trajectories, make their advantages stale, and move many ratios into the clipped
region where they supply no further improvement signal. When fresh rollout generation is
affordable, one update followed by immediate resampling is the cleaner on-policy choice.
At the end of either schedule, the system discards the short-lived buffer, copies the
updated $\pi_\theta$ into $\pi_{\mathrm{old}}$, and generates the next batch with the
improved model. Asynchronous execution adds another source of mismatch because rollout
workers may lag behind the learner, but such lag is not the only reason the two policy
versions can differ.

\paragraph{Why this is not DQN experience replay.}
Both designs store experience and can reuse it, but their admissible data ages and
purposes differ. DQN keeps a long-lived replay memory containing transitions from many
behavior-policy versions. Off-policy Q-learning permits broad reuse, while randomized
sampling reduces the strong temporal correlation between consecutive environment
transitions \citep{mnih2015dqn}. A PPO or GRPO rollout buffer is normally restricted to
the current or very recent behavior policy and discarded after limited optimization. Its
purpose is to amortize expensive data collection while remaining approximately on-policy,
not to construct a persistent historical memory or directly prevent catastrophic
forgetting. Calling both objects ``a replay buffer'' hides the policy-freshness constraint
that motivates importance ratios and clipping in PPO-style training.
